\section{おわりに}
本稿では，OWLv2の物体らしさ，Alpha-CLIP，\rev{WordNetとVGを組み合わせた語彙の木}を組み合わせ，再学習を必要とせずにカテゴリの粒度を変更できる物体アノテーションのモデルを示し，検出器と分類器を選ぶ理由を実験により示した．
COCO 2017 Validationの評価では，OWLv2の物体らしさは1枚あたり50個のボックスでCOCOの物体の87.5\%とCOCO外の物体の73.4\%を囲み，COCO外の物体の再現率が0.1程度にとどまるYOLO11xとDETRに対して，アノテーションできる物体の範囲が広いことを示した．
一方，COCOの80クラスによる$ F $は，YOLO11xのボックスを用いた構成が0.535，DETRが0.486，OWLv2が0.411であり，COCOの物体に限ればCOCOで学習した検出器が有利であった．
分類器は，全ての検出器でAlpha-CLIP（crop）がCLIPより$ F $を0.03〜0.05高め，COCOのクラスより細かい語を正しく出力した割合をCLIPの約2倍とした．
また，語彙の木の部分木ごとに確率を合計して出力の深さを決めることで，GTのボックスの階層的F値がCOCOのクラスのみの0.266から0.639に向上した．
今後の課題として，COCO外の物体を残す棄却のスコアの設計，LVISの注釈による人物以外の物体の粒度の評価，YOLO11xとOWLv2のボックスの組み合わせ，処理時間の短縮が挙げられる．
